\documentclass[aspectratio=169]{beamer}

% Theme sérieux et académique
\usetheme{Madrid}
\usecolortheme{default}

% Packages utiles
\usepackage{amsmath,amsfonts}
\usepackage{graphicx}

% Informations
\title{Time Series}
\subtitle{Basic Statistical Concepts}
\author{Aymen Ben Brik}
\institute{
Classe : 1GAMMA
}

\date{\today}

\begin{document}

% Slide titre
\begin{frame}
\titlepage
\end{frame}

\begin{frame}{Session Objectives}

By the end of this session, you will be able to:

\begin{itemize}

\item Understand the fundamental statistical concepts used in time series analysis

\item Define and interpret the moments of a distribution (mean, variance, skewness, kurtosis)

\item Understand the main estimation methods:
\begin{itemize}
\item Method of Moments (MOM)
\item Maximum Likelihood Estimation (MLE)
\end{itemize}

\item Compute and interpret classical statistical estimators

\item Understand the concept of sampling distributions

\item Introduce the basic concepts of statistical inference:
\begin{itemize}
\item Hypothesis testing
\item Confidence intervals
\end{itemize}

\end{itemize}

\end{frame}

% Slide outline
\begin{frame}{Basic Statistical Concepts}
\begin{itemize}
\item Moments of a Distribution: Fully characterizing the distribution
\item Estimation Methods: Method of Moments versus Maximum Likelihood Estimation
\item Basic Estimators: Approach and Sampling Distribution
\item Multivariate Data: Joint, marginal, and conditional distribution
\item Statistical Inference: Confidence intervals and Hypothesis Testing
\end{itemize}
\end{frame}

% Moments
\begin{frame}{Moments of a Distribution}
Moments of a random variable $X$ with density $f(x)$:

\begin{itemize}
\item $l$-th moment:
\[
m'_l = E[X^l] = \int_{-\infty}^{\infty} x^l f(x)\,dx
\]

\item $l$-th central moment:
\[
m_l = E[(X-\mu)^l] = \int_{-\infty}^{\infty} (x-\mu)^l f(x)\,dx
\]

\item Examples:
\begin{itemize}
\item Expectation: $E[X]$
\item Variance: $E[(X-\mu)^2]$
\item Skewness: $S(X) = E\left[\frac{(X-\mu)^3}{\sigma^3}\right]$
\item Kurtosis: $K(X) = E\left[\frac{(X-\mu)^4}{\sigma^4}\right]$
\end{itemize}

\end{itemize}
\end{frame}

% Estimation
\begin{frame}{Statistical Estimation}

Parametric Statistics:

Observe $x_1, \dots, x_n$ from random variables

\[
X_1, \dots, X_n \sim f(x;\theta)
\]

where $\theta$ is unknown.

\vspace{0.5cm}

Goal: Estimate $\theta$ using observations.

\vspace{0.5cm}

Approaches:

\begin{itemize}
\item Method of Moments (MOM)
\item Maximum Likelihood Estimation (MLE)
\end{itemize}

\end{frame}

% Classic estimators
\begin{frame}{Examples of Classic Estimators}

Given data $\{X_1,\dots,X_n\}$:

\begin{itemize}

\item Sample mean:

\[
\hat{\mu} = \bar{X} = \frac{1}{n}\sum_{i=1}^{n}X_i
\]

\item Sample variance:

\[
\hat{\sigma}^2 =
\frac{1}{n-1}\sum_{i=1}^{n}(X_i-\hat{\mu})^2
\]

\end{itemize}

\end{frame}
\begin{frame}{Examples of Classic Estimators 2}

Given data $\{X_1,\dots,X_n\}$:

\begin{itemize}

\item Sample skewness:

\[
\hat{S} =
\frac{1}{(n-1)\hat{\sigma}^3}
\sum_{i=1}^{n}(X_i-\hat{\mu})^3
\]

\item Sample kurtosis:

\[
\hat{K} =
\frac{1}{(n-1)\hat{\sigma}^4}
\sum_{i=1}^{n}(X_i-\hat{\mu})^4
\]

\end{itemize}

\end{frame}

% Sampling distribution
\begin{frame}{Sampling Distribution of Estimators}

Under normality assumption:

\[
X_i \sim \mathcal{N}(\mu,\sigma^2)
\]

\[
\bar{X} \sim \mathcal{N}(\mu,\sigma^2/n)
\]

\[
\frac{(n-1)S^2}{\sigma^2} \sim \chi^2_{n-1}
\]

Properties:

\begin{itemize}
\item Unbiasedness
\item Consistency
\end{itemize}

\end{frame}

% MOM
\begin{frame}{Method of Moments Estimation}

Principle:

\[
E[X^p] =
\frac{1}{n}
\sum_{i=1}^{n} X_i^p
\]

Example:

\[
X_i \sim \mathcal{N}(\mu,\sigma^2)
\]

Mean estimator:

\[
\hat{\mu} =
\frac{1}{n}\sum X_i
\]

Variance estimator:

\[
\hat{\sigma}^2 =
\frac{1}{n}\sum (X_i-\bar{X})^2
\]

\end{frame}

% MLE
\begin{frame}{Maximum Likelihood Estimation}

Likelihood function:

\[
L(\theta) =
f(x_1,\dots,x_n;\theta)
\]

Under independence:

\[
L(\theta) =
\prod_{i=1}^{n} f(x_i;\theta)
\]

MLE estimator:

\[
\hat{\theta} =
\arg\max_\theta L(\theta)
\]

\end{frame}

% Multivariate
\begin{frame}{Multivariate Distribution}

Joint distribution:

\[
f(x,y) = f(x|y)f(y)
\]

For $n$ variables:

\[
f(x_1,\dots,x_n)
=
\prod_{i=1}^{n}
f(x_i|x_{i-1},\dots,x_1)
\]

Example normal:

\[
f(x)
=
\prod_{i=1}^{n}
\frac{1}{\sqrt{2\pi\sigma_i^2}}
\exp
\left(
-\frac{(x_i-\mu_i)^2}{2\sigma_i^2}
\right)
\]

\end{frame}

% Inference
\begin{frame}{Statistical Inference}

Hypothesis testing:

\[
H_0:\theta=\theta_0
\quad vs \quad
H_a:\theta\neq\theta_0
\]

P-value:

Measure of plausibility of $H_0$

Significance level:

Probability of type I error

Confidence interval:

\[
\hat{\theta} \pm k
\]

such that

\[
P(\theta \in [\hat{\theta}-k,\hat{\theta}+k]) = 1-\alpha
\]

\end{frame}

\end{document}